\chapter*{General introduction}
\addcontentsline{toc}{chapter}{General introduction}
Artificial Intelligence has advanced to a point where failing to leverage it — even at its most basic — risks falling behind. The time saved through AI and LLMs is too valuable to ignore, allowing research \& development to be further accelerated, but also complex data-driven operations — where decisions once required teams of analysts and hours of querying — to be handled in seconds through natural conversation. This applies especially in the telecommunications field, where the industry is too data sensitive, as it generates a vast amount of network performance metrics, subscriber behavior data and commercial indicators in real time. Usually this process is done through dedicated teams, specialized tools and significant time investment. However with the emergence of LLMs and their advanced reasoning, as well as agentic AI, it is now possible for one operator to simply ask a question and receive a well informed, data-backed response promptly. The project explores this possibility further, building an AI-powered network analytics assistant that allows network operators to query, analyze and act on both network and commercial data through natural language alone in order to help improve customer experience and boost revenue.

Today, a commercial manager who would want to identify upsell opportunities in a specific region, for instance 5G upsell, and how to approach this upsell campaign, would need to file a request to a data analyst in order to consult strategies and best in slot options. Accessing the data is rarely the bottleneck. Modern telecom operators have dashboards, BI tools, etc... The real issue is what comes after extracting said data, yes it would take time to get the right query in order to get exactly what the manager wants, but what takes even more time is making sense of it: raw metrics in isolation mean little; their value emerges only when contextualized/correlated, and interpreted across multiple dimensions simultaneously. A single anomaly in network performance could be insignificant on its own, might even just be an outlier — but viewed alongside subscriber complaint patterns, billing trends and geographic coverage gaps, that same anomaly can turn out to be the first sign of something worth acting on. Arriving at the right conclusion requires not just data access, but cross-domain expertise, time and collaborative analysis across multiple teams.

Stated formally: given a telecom operator's OSS and BSS databases $D_{net}$ and $D_{com}$, and a natural-language question $q$ posed by a non-technical user, the problem is to produce an answer $A(q)$ that is correct with respect to the current state of those databases, delivered in interactive time, and where any action with operational consequence is proposed for human approval rather than executed autonomously.

In response to this challenge, this project sets out to develop an AI-powered conversational assistant for telecom operators, capable of running queries quickly, and more importantly, correctly analyze and come up with solutions/suggestions/overview, reducing dependency on technical teams for day-to-day analytical tasks.
The specific objectives, each stated with the criterion used to judge it:
\begin{itemize}
    \item \textbf{A representative simulated environment.} Two databases split along
    the OSS/BSS line, schema derived from TM~Forum's Information Framework, populated
    with $\sim$50{,}000 internally consistent subscribers and joinable on a single key.
    \item \textbf{Multi-step natural-language reasoning across both domains.} The agent
    chains more than one query when a single one cannot answer the question, and
    inspects the real schema rather than guessing column names. Target: at least 80\%
    correct on a golden set of factual questions with independently computed answers.
    \item \textbf{Responses grounded in domain knowledge.} Every query is written
    against the live schema through a knowledge graph and retrieval over worked
    examples; answers citing numbers absent from the query results are rejected and
    regenerated automatically.
    \item \textbf{Operations beyond answering.} Charts, an interactive drill-down mind
    map, a segment-to-strategy diagram, and CSV export available on eligible answers.
    \item \textbf{A continuously live environment.} A simulator advances both databases
    on a fixed tick rather than regenerating on demand, so no answer can be reused
    without re-querying. Target: no data series more than one reporting cycle stale.
    \item \textbf{Human approval for operational actions.} Every action with external
    consequence is proposed, never executed. Target: zero autonomous executions on
    adversarial prompts.
\end{itemize}

Chapter~\ref{ch:eval} reports the measured outcome against each of these criteria,
including the ones that were not met.

\subsection*{Contributions}

The specific contributions of this work are:
\begin{enumerate}
    \item A dual-database synthetic telecom environment with a \emph{causally coupled}
    simulator, in which a raised fault degrades its own cell's counters, propagates to
    the customers that cell serves, and raises their churn risk until it heals ---
    rather than each of those being an independent random process.
    \item An agent architecture that combines structured tool-calling with
    schema-graph retrieval and a deterministic post-hoc verification layer, so that
    factual correctness is enforced by code rather than requested in a prompt.
    \item A disambiguation mechanism that declines to answer an under-specified
    question, offering concrete readings of it instead of silently choosing one.
    \item A human-in-the-loop action pipeline in which the agent may propose an
    operational action but never execute it.
    \item An empirical observation that the failure modes of this system persist
    across a substantial increase in model size, locating the remaining errors in the
    retrieval and grounding layer rather than in model capacity --- with a measured
    evaluation, including failures, reported rather than omitted.
\end{enumerate}


The project follows an empirical, iterative methodology — driven by experimentation and progressive refinement rather than a predefined specification. Approaches were selected, tested, and either validated or discarded based on observed results, making the development process as much a research exercise as an engineering one.

As a result, the scope was redefined to what is feasible rather than what was initially planned; it is now bounded to a simulated telecom environment representing a hypothetical operator with around 50{,}000 subscribers across a fictional geography, covering the network domain in sites, cells, KPIs and alarms, as well as the commercial domain via subscribers, billing, churn and campaigns. Real-time hardware integration, production infrastructure, and multi-operator scenarios fall outside the current scope for now.

The produced deliverables are:
\begin{itemize}
    \item Dual database simulated telecom environment, with dynamic data simulator.
    \item A conversational AI agent capable of natural language querying and operational action execution such as excel export.
    \item A churn prediction model integrated into the live environment
    \item Real time dashboard surfacing network and commercial intelligence
    \item A RAG knowledge base grounding the agent in telecom domain
\end{itemize}

The rest of this report follows the order the work actually happened, rather than
a tidy after-the-fact architecture.

\textbf{Chapter~\ref{ch:context}} covers the host organization and the context the
project sits in --- what SmartCare is, and why building against it directly was
never an option.

\textbf{Chapter~\ref{ch:related}} positions the work against existing research in
text-to-SQL, language-model agents, and retrieval-based grounding, and is explicit
about which parts of the design are established technique and which are not.

\textbf{Chapter~\ref{chapter:BF-DFP}} lays the foundation: two synthetic databases
modelled on TM Forum's abstractions, a first agent loop running on a local model
to prove the idea works, and then the move to cloud inference once local hardware
turned out to be the ceiling rather than the idea itself.

\textbf{Chapter~\ref{ch:livingdb}} turns that static data into a living environment
--- a simulator that keeps KPIs, alarms, incidents and churn moving, so the
assistant is always reading a target that never sits still.

\textbf{Chapter~\ref{ch:agent}} covers the agent itself: the reasoning loop, how it
grounds its queries in the real schema instead of guessing, how it catches its own
mistakes before they reach the user, and everything it can do beyond answering in
plain text.

\textbf{Chapter~\ref{ch:eval}} measures all of it against a golden set, reports
where it fails, and documents a churn model that was found not to transfer.

\textbf{Chapter~\ref{ch:limits}} sets out what the system does not do, what that
costs, and what would have to happen next.
